\section{Formal Framework}
\label{sec:framework}

\subsection{System model and output abstraction}

\begin{definition}[System under test]
\label{def:sut}
A system under test (SUT) is a map $f: X \to \PY$ from an input space $X$ to a
distribution over a structured output space $Y = Y_1 \times \cdots \times Y_q$.
Deterministic systems are the degenerate case in which $f(x)$ places all mass on
a single output vector $y = (y_1,\dots,y_q)$.
\end{definition}

The components $Y_j$ are heterogeneous behavioural channels: a predicted class, a
confidence score, a robustness verdict, a measurement-outcome statistic, and so
on. Raw outputs are continuous or high-cardinality, so we abstract each channel
into a finite alphabet of \emph{behaviour symbols}.

\begin{definition}[Semantic partition and abstract output]
\label{def:partition}
A partitioner $\pi_j: Y_j \to W_j$ maps a raw output component to a symbol in a
finite alphabet $W_j$. The abstraction map is
\[
\fhat(x) \;=\; \bigl(\pi_1(y_1),\,\dots,\,\pi_q(y_q)\bigr) \in \Wspace,
\quad y \sim f(x),
\]
evaluated at the most-probable output for stochastic $f$. We write $v = \max_j
|W_j|$ for the maximum alphabet size and $v_j = |W_j|$.
\end{definition}

\subsection{Output covering arrays}

Let $[q]=\{1,\dots,q\}$. An \emph{$s$-way output tuple} fixes symbols on an
$s$-subset of channels.

\begin{definition}[Output covering array]
\label{def:oca}
An Output Covering Array $\OCA(M;s,q,w)$ is a multiset of $M$ abstract-output rows
$\{r_1,\dots,r_M\} \subseteq \Wspace$ such that for every $s$-subset
$\{j_1<\cdots<j_s\}\subseteq[q]$ and every \emph{feasible} symbol combination
$(a_1,\dots,a_s)\in W_{j_1}\times\cdots\times W_{j_s}$, some row $r_i$ satisfies
$\pi_{j_k}$-projection $r_i[j_k]=a_k$ for all $k$. A symbol combination is
\emph{feasible} if there exists an input $x\in X$ with $\fhat(x)$ exhibiting it.
\end{definition}

\begin{definition}[Output coverage and tuple efficiency]
\label{def:metrics}
For a test suite whose realised abstract outputs are $R$, the $s$-way output
coverage is
\[
\OCov_s \;=\;
\frac{\bigl|\{\text{feasible $s$-way tuples covered by } R\}\bigr|}
     {\bigl|\{\text{feasible $s$-way tuples}\}\bigr|}\in[0,1],
\]
and the tuple efficiency $\etm_s$ is the number of distinct $s$-way tuples covered
per test.
\end{definition}

\subsection{Coverage bounds}

Let $T_s$ denote the number of feasible $s$-way output tuples. Because each row
projects to exactly one tuple per $s$-subset, a single row covers at most
$\binom{q}{s}$ tuples.

\begin{theorem}[Universe size]
\label{thm:universe}
The number of $s$-way output tuples over $\Wspace$ (before feasibility pruning) is
$\sum_{\{j_1<\cdots<j_s\}} \prod_{k=1}^{s} v_{j_k}$, which for the homogeneous case
$v_j=v$ equals $\binom{q}{s}\,v^{s}$.
\end{theorem}

\begin{theorem}[Lower bound on rows]
\label{thm:lb}
Any $\OCA(M;s,q,w)$ satisfies $M \ge \max_{\{j_1<\cdots<j_s\}}\prod_{k=1}^{s}
v_{j_k}$ whenever all corresponding symbol combinations are feasible. In
particular, no single $s$-subset's tuples can be covered by fewer rows than the
number of distinct symbol combinations on that subset.
\end{theorem}

\begin{proof}
Fix the $s$-subset $J^{*}$ attaining the maximum product. Its feasible tuples are
pairwise distinct on $J^{*}$, and each row of the array contributes exactly one
symbol combination on $J^{*}$. Covering all $\prod_{k}v_{j_k}$ combinations
therefore requires at least that many rows.
\end{proof}

\begin{corollary}[$M \ge v^{s}$]
\label{cor:lb}
For the homogeneous case $v_j=v$ with all combinations feasible on some
$s$-subset, every $\OCA(M;s,q,w)$ satisfies $M \ge v^{s}$.
\end{corollary}

\begin{theorem}[Existence]
\label{thm:existence}
For any finite $\Wspace$ and strength $s\le q$, an $\OCA(M;s,q,w)$ covering all
feasible $s$-way tuples exists. Enumerating one row per feasible abstract output
is such an array; greedy construction (Algorithm~\ref{alg:rnwise}, Stage~2) yields
one of size $O(v^{s}\log \binom{q}{s})$ in the homogeneous unconstrained case,
matching the classical covering-array growth rate.
\end{theorem}

\subsection{The \rnwise{} algorithm}

Given a target abstract output $w^{*}$, inverse mapping seeks an input realising
it by minimising the mismatch loss
\begin{equation}
\label{eq:loss}
L(x;w^{*}) \;=\; \sum_{j=1}^{q} \delta\!\bigl(\pi_j(f_j(x)),\, w^{*}_j\bigr)
\;+\; \lambda\, R(x),
\end{equation}
where $\delta$ is the $0/1$ symbol mismatch, $R$ is an optional input regulariser
keeping $x$ realistic, and $\lambda \ge 0$ trades feasibility against realism. A
loss of zero means $\fhat(x)=w^{*}$ exactly.

\begin{algorithm}[t]
\caption{\rnwise{} Output-Oriented Testing}
\label{alg:rnwise}
\begin{algorithmic}[1]
\REQUIRE SUT $f$, partition scheme $\{\pi_j\}$, strength $s$, probe inputs $P$,
optimiser $\mathcal{O}$, budget \textsc{max\_iter}
\ENSURE prioritised test suite $\mathcal{S}$
\STATE \textbf{Stage 1 (Feasibility).} $\;\mathcal{F}\leftarrow\{\fhat(x): x\in
P\}$; record an exemplar input per abstract output; induce feasible $s$-way tuples.
\STATE \textbf{Stage 2 (Covering array).} build $\OCA(M;s,q,w)$ over $\mathcal{F}$
(greedy: repeatedly add the feasible output covering the most uncovered feasible
tuples; verify $M \ge v^{s}$).
\STATE \textbf{Stage 3 (Inverse mapping).} \FORALL{rows $w^{*}$ of the array}
  \STATE warm-start from the exemplar for $w^{*}$ if available;
  \STATE else $x \leftarrow \arg\min_x L(x;w^{*})$ via $\mathcal{O}$ (up to
  \textsc{max\_iter}); record realised $\fhat(x)$.
\ENDFOR
\STATE \textbf{Stage 4 (Prioritise).} greedily order tests by marginal
$\OCov_s$; truncate at the coverage plateau.
\RETURN $\mathcal{S}$
\end{algorithmic}
\end{algorithm}

The four stages are independent of the SUT's internals: Stage~1 needs only
black-box evaluations, Stage~2 is pure combinatorics over $\Wspace$, Stage~3 uses
gradient-free optimisation (so it applies to non-differentiable and quantum
systems), and Stage~4 is standard coverage-greedy prioritisation. This modularity
is what makes the same algorithm applicable across all four evaluated domains.
